\documentclass[tikz,border=2mm]{standalone}
\usetikzlibrary{math}

\makeatletter
% 1. 定义一个全局布尔开关，用于记录是否发现数字
\newif\iffounddigit

% 2. 核心检查宏：确保待查字符和目标字符串的 Catcode 一致
\newcommand{\checkfornumbers}[1]{%
  \global\founddigitfalse
  % 将输入的参数（如 \P{1}）彻底转为纯字符串，存入 \temp@target
  \edef\temp@target{\detokenize{#1}}%
  % 定义内部循环处理逻辑
  \def\do##1{%
    % 同样将 0-9 的数字转为 detokenize 格式，确保比对时 catcode 一致
    \edef\temp@d{\detokenize{##1}}%
    \expandafter\pgfutil@in@\expandafter{\temp@d}{\temp@target}%
    \ifpgfutil@in@\global\founddigittrue\fi
  }%
  % 手动展开循环，避免 \foreach 产生子组(group)导致变量无法传出
  \do0\do1\do2\do3\do4\do5\do6\do7\do8\do9
}

% 3. 定义 TikZ 键值逻辑
\tikzset{
  % 逻辑 A：处理点（例如 \P）
  define/@circle-point/.code args={#1,#2}{%
    \message{^^J[识别结果：点模式] 中心: \detokenize{#1}, 目标点: \detokenize{#2}^^J}%
    % 这里可以继续写你的 TikZ 绘图逻辑
  },
  % 逻辑 B：处理表达式或数值（例如 \P{1} 或 5cm）
  define/@circle-radius/.code args={#1,#2}{%
    \message{^^J[识别结果：数学/半径模式] 中心: \detokenize{#1}, 表达式: \detokenize{#2}^^J}%
    % 这里可以继续写你的 TikZ 绘图逻辑
  },
  % 主入口：根据 #2 是否含数字进行分发
  define/circle/.code args={#1,#2}{%
    \checkfornumbers{#2}%
    \iffounddigit
      % 如果含数字，跳转到半径/表达式处理逻辑
      \pgfkeys{/tikz/define/@circle-radius={#1,#2}}%
    \else
      % 如果不含数字，跳转到点处理逻辑
      \pgfkeys{/tikz/define/@circle-point={#1,#2}}%
    \fi
  },
}
\makeatother

\begin{document}
\begin{tikzpicture}
  \tikzmath{
    % 定义圆心分量
    \O{1} = 1; \O{2} = 1;
    % 定义圆周上的点分量
    \P{1} = 2; \P{2} = 3;
    % 定义半径变量
    \radius = 5;
  }

  % --- 测试场景 ---

  % 调用 1：识别为点（\P 不含数字字符）
  \tikzset{define/circle={\O,\P}}

  % 调用 2：识别为数学表达式（含有数字 1）
  \tikzset{define/circle={\O,\P{1}}}

  % 调用 3：识别为具体数值
  \tikzset{define/circle={\O,5cm}}

  % 调用 4：复杂表达式（含有数字 2）
  \tikzset{define/circle={\O,\P{1}+\P{2}+\O{2}^2}}

\end{tikzpicture}
\end{document}

% 之所以全部识别为“目标点”，是因为你传入的参数（如 \P{1}）在进入 \detokenize 之前，由于 LaTeX 的参数传递机制，其中的数字被包裹在了大括号 {} 里，或者在展开过程中发生了变化。

% 另一个极高概率的原因是：\pgfutil@in@ 在处理 \detokenize 生成的字符串时，对空格极其敏感。\detokenize 会在每个宏（如 \P）后面自动补一个空格。

% 要彻底解决这个问题，我们需要让检查逻辑更加“暴力”一点：直接在原始 Token 流中搜索。

% 修正后的完整方案
% 这个版本舍弃了不稳定的 \pgfutil@in@，改用一个更可靠的 递归扫描器。它会逐个字符检查，不仅能看清文字，还能看穿 {} 内部的内容。